\chapter{Urine Electrolytes}

\section*{What Is This Test?}
Urine electrolyte testing measures sodium, potassium, chloride, and sometimes other electrolytes in a spot or timed urine specimen. These results help assess renal handling of salt and water and can support evaluation of electrolyte and acid-base disorders.

\section*{Alternative Names}
\begin{itemize}
  \item Urine Sodium
  \item Urine Potassium
  \item Urine Chloride
  \item Urinary Electrolytes
  \item Urine Ion Profile
\end{itemize}
\section*{Why This Test Is Ordered}
It is used in selected cases of hyponatremia, hypokalemia, metabolic alkalosis or acidosis, suspected renal salt wasting, diuretic use, adrenal disorders, and assessment of kidney tubular function. It is not usually needed for uncomplicated electrolyte abnormalities.

\section*{How This Test Is Performed}
A spot urine sample is most common; a timed collection may be used when total daily excretion is needed. Blood electrolytes, creatinine, osmolality, and medication history are interpreted alongside the urine result.

\section*{Preparation, Risks, and Limitations}
No fasting is usually required. The test is noninvasive. Diuretics, intravenous fluids, diet, vomiting, diarrhea, kidney function, and timing relative to the electrolyte disturbance can change results. A spot urine concentration is not the same as total daily excretion.

\section*{Understanding Results}
Urine sodium and chloride help distinguish renal from extrarenal losses in selected clinical settings. Urine potassium helps assess inappropriate renal potassium loss. Interpretation requires the patient's volume status, serum values, acid-base status, and current medicines.

\section*{References}
Kidney Disease: Improving Global Outcomes guidance on evaluation of sodium and potassium disorders.

\clearpage
\section*{Understanding Results and Follow-up}
The laboratory report should identify the specimen, method, units, reference interval or decision limit, and whether the result is positive, negative, abnormal, or inconclusive. Reference intervals vary with age, sex, pregnancy, specimen type, assay, and laboratory; a result outside a range is not automatically a diagnosis, and a result within a range does not always exclude disease. Discuss unexpected results with the ordering clinician, who can decide whether repeat or confirmatory testing, treatment, or referral is appropriate. Do not change medicines or delay urgent care based on an isolated result.


\section*{Regulatory Status and Authorization Date}
This chapter describes a test category, not a specific commercial product. FDA, CDSCO, EU IVDR, and MHRA approval or registration dates therefore cannot be assigned without the assay manufacturer, product name, instrument, and intended use. For laboratory-developed tests, an individual agency approval date may not exist. See the Preface for the interpretation of regulatory status.

\clearpage
